\documentclass[10pt,twocolumn]{article}
\usepackage{fontspec}
\setmainfont{Times New Roman}
\setsansfont{Helvetica Neue}
\setmonofont{Menlo}
\usepackage[letterpaper,top=0.64in,bottom=0.70in,left=0.67in,right=0.67in,columnsep=0.24in]{geometry}
\usepackage{amsmath,amssymb,amsthm,bm,microtype,booktabs,graphicx,xcolor,enumitem}
\usepackage[colorlinks=true,linkcolor=black,citecolor=black,urlcolor=blue!55!black]{hyperref}
\graphicspath{{results/}}
\setlength{\parindent}{1em}\setlength{\parskip}{0pt}
\setlength{\textfloatsep}{7pt plus 2pt minus 2pt}
\setlength{\floatsep}{6pt plus 2pt minus 2pt}
\setlength{\intextsep}{6pt plus 2pt minus 2pt}
\setlist{nosep,leftmargin=*}\renewcommand{\arraystretch}{1.08}
\newtheorem{proposition}{Proposition}[section]
\newtheorem{theorem}[proposition]{Theorem}
\newtheorem{lemma}[proposition]{Lemma}
\theoremstyle{definition}\newtheorem{remark}[proposition]{Remark}
\newcommand{\R}{\mathbb{R}}
\newcommand{\TV}{\operatorname{TV}}
\newcommand{\diver}{\operatorname{div}}
\newcommand{\grad}{\nabla}
\newcommand{\norm}[1]{\left\lVert#1\right\rVert}
\newcommand{\code}[1]{\texttt{#1}}
\newcommand{\PiD}{\Pi}
\newcommand{\Tf}{T_f}


\title{\textbf{Transport-State Fusion for Fast\\Meyer--Bregman Decomposition}}
\author{Joshuah Rainstar\\\small\texttt{joshuah.rainstar@gmail.com}}
\date{September 4, 2026}
\begin{document}
\twocolumn[
\maketitle
\begin{abstract}
Efficient cartoon--texture decomposition requires more than inexpensive
iterations: the information needed by the next coupled subproblem must
survive the current one. We present an engineering realization of
Meyer--Bregman decomposition organized around a six-field transport state.
Two primal images and four accumulated gradient fields close a warm
interleaved recurrence. Retained primal spectra, streamed reflected
divergences, and a lower-triangular spectral solve eliminate redundant work
without changing that recurrence. The construction arose from studying
transport between related measurements in an earlier superresolution
experiment; matrix-function approximation was subsequently used to advance
the retained state, rather than to define it.

We distinguish exact implementation improvements from approximate trajectory
acceleration. A finite-flow polynomial on the complete state avoids the
specific retained-carrier and contour-halo defects of an earlier scalar
surrogate. Its numerical machinery is standard Krylov projection; the
contribution is the state formulation, its fused realization, and the
measured tradeoff on this decomposition. A September audit adds denser
ordinary controls, late trajectory checkpoints, and inexpensive feasible
primal--dual certificates. Agreement with a visually accepted 64-pass image
is not equivalent to objective convergence. A current-state factor
experiment improves finite-reference accuracy, but fixed factors have
parameter-dependent counterexamples. We identify statistically determined,
certified, low-overhead step selection as an open engineering requirement.
Code, numerical oracles, timing samples, and negative controls accompany
the paper.
\end{abstract}
\par\vspace{.6em}
\noindent\textbf{Keywords:} Meyer decomposition, bounded flux, Split Bregman,
transport state, spectral fusion, finite flow, primal--dual certificate.
\par\vspace{1em}
]

\section{Introduction}
Meyer's model represents oscillation through the divergence of a bounded
vector field, rather than by a frequency cutoff \cite{meyer2001,aujol2005}.
Gilles and Osher used duality and Split Bregman to implement the associated
alternating ROF subproblems \cite{gilles2024,goldstein2009}. That model and
its Bregman interpretation are the starting point of this work.

Our engineering question is what must remain available when one coupled
subproblem becomes the next. The driving images change during the
alternation. An estimate of the image alone does not necessarily retain
the flux and projection information that makes the next update effective.
Repeated cold solves discard useful information; extrapolation of stale
subproblem updates can carry information in the wrong frame.

The development began with transform planning and a two-lattice
superresolution experiment concerned with transport between related
measurements. In the Meyer problem, this motivated retaining the coupled
transport state and reducing repeated work around it. The resulting
sequence was warm interleaving, six-field reduction, retained spectra,
paired spectral solution, and only then finite-state advance.
Transport here refers to the model's bounded-flux structure and propagation
of its coupled state. We do not claim a new optimal-transport distance or
a mass-preserving image warp.

This distinction matters for attribution. The transport construction was
the author's motivating idea. Krylov projection is a recognized numerical
tool used in its later implementation, not the claimed source of that
idea. Conversely, a different route of discovery does not make established
matrix-function algebra new.

The engineering contributions are:
\begin{enumerate}
\item A closed warm six-field recurrence and its exact spectral fusion,
including source initialization and effective two-product output.
\item A diagnosis of a scalar shortcut's failure and a finite-flow
implementation on the retained state, with independent checks.
\item A measurement framework separating elapsed work, finite-trajectory
agreement, and objective certificates, including negative results.
\end{enumerate}
The September revision adds a current-state factor experiment and explicit
requirements for dynamic factors. These are subsequent investigations, not
retroactive explanations of the August measurements.

\subsection{Relation to prior work}
The Meyer model, Chambolle projection, and Split Bregman are established
\cite{meyer2001,aujol2005,chambolle2004,goldstein2009}.
Gilles--Osher's Bregman implementation appeared as UCLA CAM Report 11-73
in 2011; its 2024 arXiv posting is later dissemination \cite{gilles2024}.
Fourier diagonalization, warm starts, and Krylov matrix-function actions
are established techniques \cite{saad1992,jia2015}. Accelerating splitting
states has close precedents in Anderson-accelerated Douglas--Rachford and
ADMM \cite{fu2020,ouyang2020}. Those are relevant comparisons, although
they do not identify the finite polynomial or exact realization here.

A changing inner driving image does not make all acceleration of a fully
specified enlarged fixed-point map impossible. It does invalidate treating
successive inner problems as one unchanged problem without justification.
Historical experiments motivating this work rejected stale subproblem
momentum. The present method reconstructs the current coupled update.
\section{The retained transport state}

Let $X=\R^{H\times W}$ with periodic forward gradient $D:X\to X^2$ and
$D^*$ its adjoint. We use $\operatorname{div}=-D^*$ throughout. The discrete isotropic variation is
\begin{equation}
J(u)=\TV(u)=\sum_x \norm{(Du)(x)}_2 .
\end{equation}
Meyer's discrete oscillation norm is
\begin{equation}
\norm v_G=\inf_{v=D^*p}\norm p_{\infty,2},\qquad
\norm p_{\infty,2}=\max_x\norm{p(x)}_2 .
\label{eq:gnorm}
\end{equation}
The associated BV--$G$ formulations and alternating projections are
developed in \cite{meyer2001,aujol2005,gillesmeyer2024}.  We use the same
$[0,255]$ intensity convention and fixed parameters $\lambda=.05$, $\mu=40$.

The optimized reference retains two primal fields and four accumulated
gradient fields:
\begin{equation}
z=(u,w,t_u^x,t_u^y,t_w^x,t_w^y).
\label{eq:state}
\end{equation}
Here $w=f-u-v$ is the finite-pass survivor.  Consequently the public
two-product decomposition is
\begin{equation}
v=f-u-w,\qquad \bar u=f-v=u+w.
\label{eq:effective}
\end{equation}
This convention matters in comparisons: raw $u$ from a finite model is not
the effective cartoon $f-v$.

Define the pointwise disk projection and reflected projection
\begin{equation}
\PiD_a(t)=\begin{cases}t,&|t|\le a,\\a,t/|t|,&|t|>a,
\end{cases}\qquad R_a=I-2\PiD_a,
\label{eq:projection}
\end{equation}
and the screened inverse $S(c,\eta)=(cI+\eta D^*D)^{-1}$.  With
\begin{align}
c_u&=\lambda,&\eta_u&=2\lambda,&a_u&=\eta_u^{-1},\nonumber\\
c_w&=\mu^{-1},&\eta_w&=10\mu^{-1},&a_w&=\eta_w^{-1},
\end{align}
Write $S_u=S(c_u,\eta_u)$ and $S_w=S(c_w,\eta_w)$.
One exact warm fused pass $\Tf:z\mapsto z^+$ is
\begin{align}
u^+&=S(c_u,\eta_u)
 \{c_u(u+w)+\eta_uD^*R_{a_u}(t_u)\},\label{eq:upass}\\
w^+&=S(c_w,\eta_w)
 \{c_w(f-u^+)+\eta_wD^*R_{a_w}(t_w)\},\label{eq:wpass}\\
t_u^+&=Du^++\PiD_{a_u}(t_u),\qquad
t_w^+=Dw^++\PiD_{a_w}(t_w).\label{eq:tpass}
\end{align}

\begin{proposition}[Lower-triangular spectral pass]
For periodic boundaries, equations \eqref{eq:upass}--\eqref{eq:wpass} form a
lower-triangular two-by-two system at every Fourier coefficient.  Both
reflected divergences may be transformed before either inverse transform,
and the two primal spectra may be solved in one range traversal.
\end{proposition}
\begin{proof}
$D^*D$ diagonalizes under the periodic DFT.  If $\sigma_u$ and $\sigma_w$
are the two screened multipliers and $\widehat d_i$ denotes the transform
of $D^*R_{a_i}(t_i)$, then
\begin{align*}
\widehat u^+&=\sigma_u\{c_u(\widehat u+\widehat w)
+\eta_u\widehat d_u\},\\
\widehat w^+&=\sigma_w\{c_w(\widehat f-\widehat u^+)
+\eta_w\widehat d_w\}.
\end{align*}
The second equation depends on the newly computed first variable and no
reverse coupling remains.  This is exactly the paired spectral loop in the
C++ kernel.
\end{proof}


\subsection{Source initialization}
The first pass is source driven:
\begin{align}
u_1&=S_u(c_uf),&w_1&=S_w\{c_w(f-u_1)\},\nonumber\\
t_{u,1}&=Du_1,&t_{w,1}&=Dw_1.
\label{eq:first}
\end{align}
It is not obtained by inserting a zero state into \eqref{eq:upass}.
Every experimental pass count includes this initialization.
The plus signs before $D^*R$ in the recurrence correspond to the code's
$-\eta\operatorname{div}R$. Exact recomposition does not itself certify
that the finite emitted texture lies in the prescribed $G$ ball.

\section{Exact implementation economy}
The current driving images are $u+w$ and $f-u^+$. They change, but the
fields needed to compute them and their Bregman corrections remain in the
state. Warm interleaving takes one inner sweep per route and immediately
updates the coupled state. It is a different finite schedule from converging
both inner problems at every outer step.

The screened solve resolves global spatial linear coupling exactly for the
selected boundary conditions. This can replace many local updates even
when each local update is individually cheap. The retained state and
global solve work together; neither a warm start nor a cheap local step
alone explains the efficiency.

Retaining $\widehat f,\widehat u,\widehat w$ avoids transforming driving
images already known spectrally. Streaming the reflected divergence avoids
materializing unneeded intermediate fields. Pairing the spectral solution
shares traversal and synchronization work. An ordinary outer pass contains
two forward and two inverse two-dimensional transforms, one pair per
screened subproblem.

These are exact changes of representation and execution of the chosen
recurrence. They must be distinguished from taking fewer nonlinear passes
or approximating several passes by a jump. Persistent workers and reusable
plans support the native implementation. The one-axis FACR route extends
shape support, but the finite-flow tangent remains full spectral.
\section{Diagnosis of the scalar shortcut}

The rejected hard construction replaced the state recurrence by two scalar
observations.  Let $\mathcal S$ be its scalar separator, $H$ its contraction,
$P_K=I-H^K$, and
\begin{align}
s_0&=\mathcal S(f),&m_0&=P_K(f-s_0),\nonumber\\
r_1&=f-m_0,&s_1&=\mathcal S(r_1).
\end{align}
Before its final $G$-capacity correction, the texture proposal is
$m=P_K(f-s_1)$.  Write $v_H=m+q_\mu$, where $q_\mu$ is the change caused by
that final capacity route.

\begin{proposition}[Common cause of halo and retained texture]
Let $(u_*,v_*)$ be any exact two-product reference, $f=u_*+v_*$.  The second
hard observation and its final texture error satisfy
\begin{align}
r_1&=u_*+H^Kv_*+P_K(s_0-u_*),\label{eq:secondobs}\\
v_H-v_*&=P_K(u_*-s_1)-H^Kv_*+q_\mu .\label{eq:defect}
\end{align}
\end{proposition}
\begin{proof}
Since $P_K=I-H^K$,
\begin{align*}
r_1&=f-P_K(f-s_0)=s_0+H^K(f-s_0)\\
&=u_*+H^Kv_*+P_K(s_0-u_*),
\end{align*}
which proves \eqref{eq:secondobs}.  Similarly,
\begin{align*}
m-v_*&=P_K(f-s_1)-v_*\\
&=P_K(u_*-s_1)+(P_K-I)v_*\\
&=P_K(u_*-s_1)-H^Kv_*.
\end{align*}
Adding $q_\mu=v_H-m$ proves \eqref{eq:defect}.
\end{proof}

The two visible failures are therefore one representation error.
$H^Kv_*$ is the uncompleted carrier tail; it remains in cartoon with the
opposite sign in texture.  If $e=u_*-s_1$ contains a step mismatch and
$H^Ke=h_K*e$ for a normalized smoothing kernel $h_K$, then
$P_Ke=e-h_K*e$ has the paired positive and negative contour response of a
halo.  It has zero mean, so a one-sided clamp cannot remove it faithfully.
For a Fourier carrier with contraction gain $0<\rho<1$,
$H^Kv_*=\rho^Kv_*\ne0$ at every finite $K$.

Figure \ref{fig:defect} renders the terms independently.  Across six
$256^2$ controls, direct substitution into \eqref{eq:secondobs} and
\eqref{eq:defect} leaves at most $1.42\times10^{-13}$ in $L_\infty$.  This is
an identity check, not a fitted explanation.

\begin{figure*}[t]
\centering\includegraphics[width=.98\textwidth]{revision_defect.png}
\caption{Three representative controls from the six-control scalar-defect
archive. The carrier, structural, and capacity columns sum to the actual
texture error in the last column, to roundoff. Pure edge and pure carrier
isolate the terms; supported texture shows their superposition. Signed
error columns share a color scale within each row.}
\label{fig:defect}
\end{figure*}


The identity explains the observed defects of this construction. It is not
an impossibility theorem for all scalar algorithms, nor proof that a
different representation cannot produce a halo by another mechanism.
The engineering response was to advance the actual state instead of
repairing the output using a content-admission rule.

\section{Finite flow of the complete state}

After a short ordinary prefix, fix the current state $z$ and define
\begin{equation}
r=\Tf(z)-z,\qquad A\in\partial\Tf(z).
\label{eq:residual}
\end{equation}
If $\Tf$ were affine on the visited chart, its exact finite displacement after
$m$ more passes would be
\begin{equation}
\Tf^m(z)-z=p_m(A)r,\qquad
p_m(A)=I+A+\cdots+A^{m-1}.
\label{eq:finitepoly}
\end{equation}
This is the finite-trajectory target.  The stationary Newton displacement
$(I-A)^{-1}r$ was tested and rejected: it can overshoot a desired finite
state at structural transitions.

The only nonlinear branch derivative comes from the disk projection already
present in \eqref{eq:projection}.  For $n=t/|t|$,
\begin{equation}
D\PiD_a(t)h=\begin{cases}
h,&|t|<a,\\[2pt]
\dfrac{a}{|t|}(I-nn^T)h,&|t|>a,
\end{cases}
\label{eq:diskderivative}
\end{equation}
and $DR_a(t)h=h-2D\PiD_a(t)h$.  At $|t|=a$ the implementation chooses the
interior member of the Clarke derivative.  Differentiating the fused triangle
gives
\begin{align}
\dot u^+&=S_u\{c_u(\dot u+\dot w)
+\eta_uD^*DR_{a_u}(t_u)\dot t_u\},\nonumber\\
\dot w^+&=S_w\{-c_w\dot u^+
+\eta_wD^*DR_{a_w}(t_w)\dot t_w\},\label{eq:tangent}\\
\dot t_u^+&=D\dot u^++D\PiD_{a_u}(t_u)\dot t_u,\nonumber\\
\dot t_w^+&=D\dot w^++D\PiD_{a_w}(t_w)\dot t_w.\nonumber
\end{align}
No statistic of $f$ enters \eqref{eq:tangent}; only the current optimization
state and the original projection geometry do.

\subsection{A standard small numerical representation}

The implementation evaluates $r$, $Ar$, and $A^2r$.  Let
\begin{equation}
\begin{gathered}
B=[r,Ar],\qquad G=B^TB,\\
K=B^TAB,\qquad \widehat A=G^{-1}K.
\end{gathered}
\label{eq:galerkin}
\end{equation}
Because $Ar$ is the second basis vector, the first column of $\widehat A$ is
exactly $(0,1)^T$.  Only the projection of $A^2r$ needs a two-by-two solve.
The applied displacement is
\begin{equation}
\delta_m=Bp_m(\widehat A)e_1.
\label{eq:krylovjump}
\end{equation}
This is standard Krylov matrix-function machinery \cite{saad1992,jia2015}, not a newly claimed algebraic principle. It is algebraically the depth-two Arnoldi polynomial without normalization
or full-field modified Gram--Schmidt sweeps.  A fixed 64-chunk reduction
computes the five required inner products in a thread-count-independent
summation order.

\subsection{What error remains}

Write the semismooth expansion at $z$ as
\begin{equation}
\Tf(z+d)=\Tf(z)+Ad+N_z(d).
\end{equation}
Let $d_k=\Tf^k(z)-z$, $p_0=0$, $p_{k+1}=r+Ap_k$, and
$e_k=d_k-p_k$.  Then
\begin{equation}
e_{k+1}=Ae_k+N_z(d_k).
\end{equation}
If $\norm A\le\gamma$ and $\norm{N_z(d)}\le L\norm d^2/2$ on the traversed
chart, induction gives
\begin{equation}
\norm{e_m}\le \frac{L}{2}\sum_{k=0}^{m-1}
\gamma^{m-1-k}\norm{d_k}^2.
\label{eq:remainderbound}
\end{equation}
The quadratic assumption is not guaranteed across projection boundaries by
semismoothness alone. At $t=(a,0)$, select the interior derivative and take
outward $h=(\epsilon,0)$. Then
$\Pi_a(t+h)-\Pi_a(t)-h=-h$, a first-order remainder. The exterior disk
derivative also varies within an unchanged active set. Thus an active set
is not an affine chart.

The implemented error additionally contains the depth-two projection term
\begin{equation}
\norm{p_m(A)r-Bp_m(\widehat A)e_1}.
\end{equation}
Ordinary passes after each jump refresh the nonlinear disk branches and
shorten the chart over which \eqref{eq:remainderbound} accumulates.  Thus the
remaining discrepancy from pass 64 consists of finite target tail, Krylov
projection, and semismooth chart error.  The terms $P_K(u_*-s_1)$ and
$-H^Kv_*$ of \eqref{eq:defect} are absent as terms of that particular surrogate identity: this method never
forms $s_1$, $P_K$, or a scalar surrogate for $z$.

\subsection{Fixed schedules and the shorter-jump ablation}

For prefix $p$, two tangent products, $s$ ordinary refresh passes, and $j$
jumps, the exact operator-pass cost is
\begin{equation}
C=p+j(1+2+s),
\label{eq:cost}
\end{equation}
where the one counts the exact nonlinear residual evaluation.  The public
schedules are fixed independently of the source:
\begin{center}\small
\begin{tabular}{lrrrrr}\toprule
&$p$& horizon &$s$&$j$&$C$\\\midrule
fast&4&18&2&3&19\\
quality&4&10&2&5&29\\\bottomrule
\end{tabular}
\end{center}

We also tested the natural alternative ``jump less, refresh once.''  A cheap
post-jump graph retraction,
\begin{equation}
t\leftarrow Du+\PiD_a(t-Du),
\label{eq:retract}
\end{equation}
restores the dual feasibility $|t-Du|\le a$ in $O(HW)$ pointwise work and no
FFT.  Retraction alone prevents instability but loses most of the
acceleration.  Shorter jumps with one exact refresh won at $128^2$ but did
not uniformly beat the two-refresh schedules at $256^2$.  Reducing only the
horizon traded small wins among controls rather than producing dominance.
This negative ablation shows that refreshes are nonlinear chart updates, not
backtracking overhead.  A useful constrained variant is therefore likely to
cost about the same total spectral work unless it exploits a stronger bound
than \eqref{eq:retract}.

\section{C++ realization}

The public full-spectral entry point is
\path{bfft_meyer_split_flow_jump}.  Its plan owns the ordinary reduced
state and lazily allocates three six-field Krylov vectors plus primal spectra.
The implementation performs:
\begin{enumerate}[label=\arabic*.]
\item a source FFT and ordinary fused prefix;
\item one exact residual pass without overwriting live spectra;
\item two streamed tangent passes using \eqref{eq:diskderivative};
\item one fixed-chunk five-scalar Gram reduction;
\item the $2\times2$ polynomial in registers and one state AXPY;
\item fixed ordinary refresh passes, repeated by schedule.
\end{enumerate}
No candidate scan or line search is performed.  Texture is emitted as
$f-u-w$ and cartoon as $f-v$, guaranteeing exact two-product recomposition.

The Python binding exposes \path{MeyerPlan.split_flow_jump}; the default
full-spectral \code{split} uses the quality schedule.  The effect API and
still-image demo expose \code{quality}, \code{fast}, \code{fused64}, and the
old hard control.  The video demo uses the fixed fast schedule.  FACR plans
retain their configured ordinary fused recurrence because the current
finite-flow tangent is full-spectral; they also fold the finite survivor into
effective cartoon.

Native tests compare the C++ output to an independent NumPy implementation of
\eqref{eq:upass}--\eqref{eq:krylovjump}, check exact recomposition, translation
equivariance, and bit identity across thread counts, and require both public
schedules to beat their equal-cost ordinary controls on the surgical battery.


\section{Distance to the objective}
A visually accepted finite image is an engineering reference, not a
certificate of the discrete constrained model
\begin{align}
F^\star=\min_{u,\ v\in\mathcal G_\mu}
&J(u)+\frac{\lambda}{2}\norm{f-u-v}^2,\nonumber\\
&\mathcal G_\mu=\{D^*g:|g(x)|\le\mu\}.
\label{eq:objective}
\end{align}
We report $E_{64}=\norm{v_{\rm out}-v_{64}}/\norm{v_{64}}$ and late
trajectory differences, while separately recovering feasible bounds.

From any retained state define
\begin{align}
p&=\eta_u\Pi_{a_u}(t_u),&q&=D^*p,\nonumber\\
g&=(\eta_w/c_w)\Pi_{a_w}(t_w),&v_F&=D^*g.
\end{align}
The flux radii are one and $\mu$ respectively.
\begin{proposition}[Feasible objective certificate]
For
\begin{align}
P&=J(u)+\frac{\lambda}{2}\norm{f-u-v_F}^2,\nonumber\\
L&=\langle f,q\rangle-\frac{\norm q^2}{2\lambda}-\mu J(q),
\end{align}
one has $L\le F^\star\le P$.
\end{proposition}
\begin{proof}
The upper bound uses a feasible pair. For the lower bound use
$J(u)\ge\langle u,q\rangle$ and
$\sup_{v\in\mathcal G_\mu}\langle v,q\rangle=\mu J(q)$, then minimize the
quadratic. Equivalently, the gap decomposes into nonnegative terms:
\begin{align}
P-L={}&J(u)-\langle u,q\rangle+\mu J(q)-\langle v_F,q\rangle\nonumber\\
&+\frac{\lambda}{2}\norm{f-u-v_F-q/\lambda}^2.
\label{eq:gap}
\end{align}
\end{proof}
This needs pointwise projections, local differences, and reductions, with
no additional screened solve. It certifies the recovered feasible pair,
not automatically the emitted texture. We also record their RMS difference.
A small gap does not establish uniqueness of the decomposed images.

\section{Engineering evidence}
\subsection{Archived finite-flow measurements}
The August surgical battery uses six designed $256^2$ controls and
$(\lambda,\mu)=(.05,40)$ on $[0,255]$ intensities. Mean texture RMS error
to fused 64 is 3.431 for the scalar shortcut, 1.400 for fast versus 1.496
for ordinary 19, and 1.039 for quality versus 1.196 for ordinary 29.
Every equal-operator-count comparison improves in that battery.

Archived eight-thread M4 medians on a multiscale crossing are 11.52 ms
(fast), 18.16 ms (quality), and 34.60 ms (ordinary 64) at $256^2$.
These are different-accuracy operating points. Cold nested Gilles timings
in the archive use NumPy/SciPy and different stopping rules; they are
historical context, not evidence of a same-implementation speedup.

\begin{figure*}[t]
\centering
\includegraphics[width=.94\textwidth]{revision_natural.png}
\caption{Archived two-product comparison, with survivor folded into effective
cartoon. Barbara is $512^2$; the other two sources are $256^2$.
The scalar shortcut retains oscillatory material and exhibits contour
error. Finite-flow outputs track the ordinary 64-pass allocation more
closely. Texture columns share a color scale within each row.}
\end{figure*}

\subsection{September audit and clean timing rerun}
The audit adds ordinary 32, 40, and 48 controls; ordinary 2048 and 4096
checkpoints; and the certificate \eqref{eq:gap}. All four comparison
sources are resized or retained at $256^2$, including Barbara.
Parameters and periodic boundaries are unchanged.
The late reference is checked: texture drift from pass 2048 to 4096 is
0.39--1.73\%, while the relative recovered objective gap at 4096 is
$2.05\times10^{-5}$ to $1.36\times10^{-4}$.
Ordinary 64 differs from pass 4096 by 8.7--22.7\% in texture norm.

The original September screen ran under competing CPU load. Its numerical
outputs remain valid, but its timings are not used as a clean speed claim.
The new record was collected after those workers were stopped.
Plan construction and post-run diagnostics are excluded; returned-output
allocation is included. Two warmup rounds precede 15 timed rounds, with
method order shuffled in each round. Table~\ref{tab:clean} is generated
from serialized measurements.

\input{results/september_clean_table.tex}

Equal operator counts do not mean equal time. Zero error to pass 64 is
true by definition for ordinary 64, not evidence of objective convergence.
The quality schedule has a larger recovered gap than ordinary 32 on all
four sources, despite better $E_{64}$. Its role is finite-state prediction,
not direct minimization of that certificate.

\subsection{A current-state factor experiment}
The audit tests, after the four-pass prefix,
\begin{equation}
z^+=z+\alpha(\Tf(z)-z).
\label{eq:relax}
\end{equation}
This is classical full-state overrelaxation, not a new acceleration
principle. Both driving fields are recomputed inside $\Tf$ before applying
the factor. An experimental native version fuses the state combination
and maintains spectra without extra FFTs.

At $\alpha=1.75$ and 32 total passes, $E_{64}$ is 59--65\% lower than
the quality preset on the four audit sources. Recovered gaps are smaller
by factors 1.89--3.31. At 64 passes, error to pass 4096 improves by
12--27\% against ordinary 64 and the gap improves by 36--54\%, but relaxed
passes cost more. Same-count gains alone are not matched-time
convergence-speed claims.

The factor is not universal. At 128 passes, a screen of eight sources at
three parameter pairs leaves all outputs finite, yet $\alpha=1.75$
improves the gap in only 16 of 19 nontrivial cases. Counterexamples include
a ramp, an edge, and a low-frequency carrier. At $\alpha=1.9$, one ramp
has approximately 94.5 times the ordinary gap. Five other cases have
ordinary gaps below $10^{-8}$ and are excluded from gain ratios.
Two final ordinary steps improve the gap in every tested finishing-pass
ablation, but can worsen agreement with pass 64.

\section{Requirements for dynamic factors}
A constant multiplier diagnoses available progress; it is not a sufficient
method. A production factor should be determined from current transport
statistics, carry a stated mathematical guarantee, and cost materially
less than the work it saves. There are three separate requirements:
\begin{enumerate}
\item \emph{A valid statistic.} Boundary distance, nonlinear defect,
residual, or objective bound must measure the claimed quantity.
Active-set persistence alone does not certify a long affine prediction.
\item \emph{A valid decision.} A local slope can suggest a candidate.
An established bound or checked inequality must certify its advertised
property. A feasible gap comparison certifies that comparison; it does not
by itself prove convergence of repeated choices.
\item \emph{A measured cost.} Reductions, stencils, trial projections,
and rejected candidates must all be timed. Absence of extra FFTs does not
imply negligible overhead.
\end{enumerate}
The certificate \eqref{eq:gap} enables an initial investigation without
changing the model. Dynamic-factor work is maintained as a separate
follow-up. No successful adaptive convergence theorem or universal dynamic
speedup is asserted in this revision.

\section{Reproducibility and limitations}
The full-spectral implementation is checked against an independent NumPy
six-field oracle. Tests cover adjoint signs, tangent finite differences,
initialization through trajectory agreement, recomposition, translation
equivariance, and thread reproducibility. The audit adds the nonnegative
gap decomposition, a zero-gap constant case, a boundary remainder
counterexample, and native agreement for experimental factors and refreshes.

{\small\raggedright Code and evidence reside in
\path{src/detail/meyer_kernel.hpp},
\path{experiments/meyer_semismooth_state_jump.py},
\path{experiments/meyer_hard_jump_defect_proof.py}, and
\path{experiments/meyer_transport_audit/}.
Raw timing samples, arrays, counterexamples, and reproduction commands are
retained. The audit wrapper does not replace production presets.\par}

The study does not establish global convergence of warm one-sweep
interleaving for every parameter or boundary condition, image uniqueness,
or priority for every ingredient. Finite-flow evidence covers a limited
source set and primarily one parameter pair. Factors were selected during
exploration on these sources; repeated timing is not held-out validation.
The arbitrary-shape route is not included in the finite-flow speed claim.

\section{Conclusion}
The principal engineering result is a retained transport state that makes
the changing coupled problems inexpensive to reconstruct and globally
solve. Spectral fusion realizes the same recurrence with less repeated
work. Finite-flow approximation and current-state overrelaxation are
subsequent choices about how far to advance that state.

A scalar shortcut loses identifiable information; a faithful finite-state
jump improves finite-reference agreement; objective certificates reveal
progress that a finite image reference cannot measure. Further
acceleration should be earned by a cheap, formally stated, continuously
updated decision about the current transport, with its full cost included.

\section*{Research attribution}
Joshuah Rainstar developed the transport-oriented research direction,
connected it to earlier superresolution work, and supplied visual and
computational falsification criteria. Substantial algebraic exploration,
implementation, benchmark scaffolding, and manuscript drafting were
performed with AI assistance under that direction. Established model,
splitting, and matrix-function ideas are credited in their respective
sections. Claims are limited to the stated derivations and executable
evidence.
\begin{thebibliography}{10}
\small
\bibitem{meyer2001} Y. Meyer, \emph{Oscillating Patterns in Image Processing
and in Some Nonlinear Evolution Equations}. AMS, 2001.
\bibitem{aujol2005} J.-F. Aujol, G. Aubert, L. Blanc-F\'eraud, and A.
Chambolle, ``Image decomposition into a bounded variation component and an
oscillating component,'' \emph{J. Math. Imaging Vision}, 22:71--88, 2005.
\bibitem{chambolle2004} A. Chambolle, ``An algorithm for total variation
minimization and applications,'' \emph{J. Math. Imaging Vision}, 20:89--97,
2004.
\bibitem{goldstein2009} T. Goldstein and S. Osher, ``The Split Bregman method
for $L^1$-regularized problems,'' \emph{SIAM J. Imaging Sci.}, 2:323--343,
2009.
\bibitem{gilles2024} J. Gilles and S. Osher, ``Bregman implementation of
Meyer's $G$-norm for cartoon + textures decomposition,'' UCLA CAM Report 11-73, 2011; later posted as arXiv:2410.22777, 2024.
\bibitem{gillesmeyer2024} J. Gilles and Y. Meyer, ``Properties of BV--G
decomposition models,'' arXiv:2411.04456, 2024.

\bibitem{saad1992} Y. Saad, Analysis of some Krylov subspace approximations
to the matrix exponential operator. \emph{SIAM J. Numer. Anal.},
29(1):209--228, 1992.
\bibitem{jia2015} Z. Jia and H. Lv, A posteriori error estimates of Krylov
subspace approximations to matrix functions. \emph{Numerical Algorithms},
69:1--28, 2015; arXiv:1307.7219.
\bibitem{fu2020} A. Fu, J. Zhang, and S. Boyd, Anderson accelerated
Douglas--Rachford splitting. \emph{SIAM J. Sci. Comput.},
42(6):A3560--A3583, 2020.
\bibitem{ouyang2020} W. Ouyang, Y. Peng, Y. Yao, J. Zhang, and B. Deng,
Anderson acceleration for nonconvex ADMM based on Douglas--Rachford
splitting. \emph{Computer Graphics Forum}, 2020; arXiv:2006.14539.
\end{thebibliography}
\end{document}
